\documentclass[11pt]{article}
\usepackage[utf8]{inputenc}
\usepackage[a4paper,margin=1in]{geometry}
\usepackage{amsmath}
\usepackage{booktabs}
\usepackage{array}
\usepackage{enumitem}
\usepackage{microtype}
\usepackage[hidelinks]{hyperref}

\title{NLP Project --- Part 4: Analysis of the TF-IDF VSM Baseline}
\author{Soham Shah (ME22B191)}
\date{}

\begin{document}
\maketitle

\section{Setup and reproduction of baseline metrics}

The baseline uses the pipeline already implemented in Parts 2--3: Punkt
sentence segmentation, Penn-Treebank tokenization, Porter stemming, NLTK
English stop-word removal, then a TF--IDF inverted index with cosine ranking.
Running the full evaluation on the 225 Cranfield queries reproduces the
numbers cited in Table~\ref{tab:baseline}. We use these as the reference
point for the failure analysis that follows.

\begin{table}[h]
\centering
\small
\begin{tabular}{cccccc}
\toprule
$k$ & P@$k$ & R@$k$ & F$_{0.5}$@$k$ & AP@$k$ & nDCG@$k$ \\
\midrule
1  & 0.6489 & 0.1092 & 0.3034 & 0.1092 & 0.6489 \\
3  & 0.4859 & 0.2287 & 0.3672 & 0.2108 & 0.5237 \\
5  & 0.3858 & 0.2884 & 0.3388 & 0.2488 & 0.4711 \\
10 & 0.2836 & 0.4100 & 0.2892 & 0.3026 & 0.4636 \\
\bottomrule
\end{tabular}
\caption{TF--IDF VSM baseline on Cranfield (mean across 225 queries).
The full table for $k=1\ldots 10$ is in
\texttt{part4\_output/aggregate\_metrics.csv}.}
\label{tab:baseline}
\end{table}

Two observations from preprocessing diagnostics frame the rest of the
analysis. First, 46 of 225 queries (20.4\%) contain at least one token that
does not appear in the index after preprocessing. Second, two documents
(IDs 471 and 995) are empty after stop-word removal and are therefore
unreachable. Both are direct artefacts of the bag-of-words representation.

\section{Failure mode analysis}

We partition the failing queries into three structurally distinct modes.
For each mode we explain the structural cause in the VSM, give two or three
concrete query examples taken from the diagnostics in
\texttt{part4\_output/per\_query\_diagnostics.csv}, and compare against the
three Part--5 systems (BM25, LSA at $k=300$, WordNet Query Expansion) to
test whether the mode is fixable.

\subsection{Mode 1 --- Vocabulary mismatch (synonymy)}

\textbf{Cause.} Cosine similarity over a TF--IDF vocabulary scores documents
purely by surface token overlap. When a query and a relevant document
describe the same concept with different surface vocabulary --- a true
synonymy gap --- the VSM has no representation to bridge them.

\textbf{Query 12} --- ``how can the aerodynamic performance of channel flow
ground effect machines be calculated.''
(P@10 = 0.10, $n_{\text{rel}}=6$, first relevant at rank 10.) The query
phrase \emph{ground effect machines} is a 1960s term for what the Cranfield
documents call \emph{hovercraft} and \emph{peripheral jets}. The top
retrieved documents (Doc 941, Doc 1221, Doc 270) match on the generic
\emph{channel flow} tokens. The missed relevant Doc 86 (\emph{inviscid-%
incompressible flow theory of static peripheral jets}) and Doc 649
(\emph{The hovercraft --- a new concept in maritime transport}) use no
overlapping surface tokens with the query and are invisible to the VSM.

\textbf{Query 30} --- ``papers on flow visualization on slender conical
wings.'' (P@10 = 0.10, $n_{\text{rel}}=8$, first relevant at rank 6.) The
relevant papers describe the same idea using \emph{vapour screen method},
\emph{flow studies}, and \emph{flow patterns} rather than \emph{flow
visualization}. A particularly diagnostic case is Doc 633, which appears in
the top retrieved set without being relevant: it is a theoretical paper on
\emph{generalised conical flows} that shares \emph{flow}, \emph{conical},
and \emph{wings} tokens with the query but is a pure theory paper not
about visualization. This is the false-positive flip side of Mode 1 ---
surface overlap pulling in apparently relevant but actually unrelated
papers.

\textbf{Part--5 comparison.} LSA closes the synonymy gap most reliably
because it projects documents and queries into a shared latent space where
\emph{hovercraft} and \emph{ground effect machines} co-occur in training
contexts. BM25 cannot fix Mode 1 because its TF saturation does not
introduce new vocabulary mappings. WordNet QE helps only when the
synonym pair happens to be in WordNet (it is not, for the aerospace jargon
above).

\subsection{Mode 2 --- Term importance beyond TF--IDF}

\textbf{Cause.} The VSM assigns each query token exactly one weight: its
IDF. Every content token contributes positively to the cosine score with
any document that contains it. There is no representation of (a)
\emph{negation}, (b) \emph{question focus} --- which token is the topic
versus scaffolding, or (c) \emph{constraint scope} --- which token is a
modifier limiting the query. IDF reweighting cannot fix this because it
operates on token identity, not on linguistic role.

\textbf{Query 22 --- negation.} ``did anyone else discover that the
turbulent skin friction is \textbf{not} over sensitive to the nature of the
variation of the viscosity with temperature.'' (P@10 = 0.00,
$n_{\text{rel}}=2$, first relevant at rank 970.) The query asserts an
\emph{insensitivity} result. The VSM treats \emph{sensitive},
\emph{viscosity}, \emph{temperature}, \emph{skin}, \emph{friction} as
positive match signals and retrieves Doc 254 and Doc 125 ---- generic
skin-friction papers, exactly the affirmative version of the query. The
two relevant documents (Doc 68 on air-helium simulation, Doc 502 on
Squire's compressibility transformation) describe the actual insensitivity
finding using entirely different surface vocabulary and are missed.

\textbf{Query 176 --- exclusion.} ``some approximate analytical heat
conduction solutions using methods \textbf{other than} Biot's principle.''
(P@10 = 0.10.) The exclusion marker \emph{other than} is invisible to the
VSM. Because \emph{Biot} has high IDF, Doc 542 \emph{Biot's variational
principle in heat conduction} --- the paper the query explicitly tries to
exclude --- is retrieved at rank 2. Doc 582 (\emph{the melting of finite
slabs}, which uses the heat-balance integral, a genuinely different method)
is missed.

\textbf{Query 151 --- focus drift.} ``what is the best theoretical method
for calculating pressure on the surface of a wing \textbf{alone}.''
(P@10 = 0.00, first relevant at rank 19.) The constraint \emph{wing alone}
(as opposed to wing-body) is the substantive content. The system locks on
to \emph{wing} (IDF = 1.87, present in all five non-relevant top
documents) and retrieves generic wing papers (Doc 433, Doc 251, Doc 678).
The relevant papers (Docs 687, 1062, 1074--1077) all share
\emph{theoretical} but the VSM weighting cannot identify \emph{theoretical
method} as the question's focus.

\textbf{Query 63 --- scaffolding.} ``where can I \textbf{find} pressure
data on surfaces of swept cylinders.'' (P@10 = 0.00.) The query-asking
verb \emph{find} (IDF = 3.72) is matched against Doc 738, \emph{Finding
zero's of arbitrary functions}, which ranks first. The VSM has no
syntactic awareness that \emph{find} is scaffolding and \emph{swept
cylinders} is the topic.

\textbf{Part--5 comparison.} On Q22, all four systems give P@10 = 0.00 (the
first-relevant rank is 970 for VSM, 1038 for BM25, 139 for LSA, 483 for
QE); LSA narrows the gap but cannot recover the polarity. On Q176, BM25
reaches P@10 = 0.30 by saturating the TF of \emph{Biot}, but still ranks
the excluded Biot paper first --- the exclusion itself is not represented.
On Q151 only BM25 manages a marginal P@10 = 0.10. On Q63 all four systems
give P@10 = 0.00. None of BM25 (re-weighted TF), LSA (latent re-projection),
or WordNet QE (synonym broadening) introduces a role-sensitive mechanism,
which is consistent with the structural argument: Mode 2 needs syntactic
or contextual modelling, not term reweighting.

\subsection{Mode 3 --- Out-of-vocabulary query tokens}

\textbf{Cause.} The VSM's IDF table is built from the corpus vocabulary.
A query token whose stem is absent from this table has no IDF entry and
is silently dropped from the query vector. When that token is the
discriminative term of the query, the residual vector is too generic to
locate the relevant documents. In the Cranfield diagnostics this affects
46 of 225 queries.

\textbf{Query 207 --- hyphenated compound.} ``how do large changes in new
mass ratio quantitatively affect \textbf{wing-flutter} boundaries.''
(P@10 = 0.00, first relevant at rank 59.) After stemming, the compound
becomes \texttt{wing-flutt}, which appears in zero documents --- the corpus
consistently writes \emph{wing flutter} as two tokens. The surviving query
\{large, change, new, mass, ratio, affect, boundary\} is generic; the
system retrieves boundary-layer papers (Doc 458, Doc 365, Doc 1185) and
all six relevant flutter papers (Doc 1290, Doc 879, etc.) are missed.

\textbf{Query 117 --- slash-bracketed phrase.} ``is there any information
on how the addition of a \textbf{/boat-tail/} affects the normal force on
the body \ldots'' (P@10 = 0.00, first relevant at rank 18.) The Cranfield
source format wraps some technical phrases in slashes; the Penn-Treebank
tokenizer treats the leading slash as part of the token. The surface form
\texttt{/boat-tail/} exists in no document body, where the term appears
unbracketed. The relevant boat-tail papers (Doc 360, 605, 896) are missed.
This is a pure preprocessing artefact: query and documents use the same
word, kept apart by a slash.

\textbf{Query 9 --- same artefact.} ``papers on internal \textbf{/slip
flow/} heat transfer studies.'' (P@10 = 0.00.) Both \texttt{/slip} and
\texttt{flow/} are effectively OOV. The query reduces to \{paper,
internal, heat, transfer, studies\}, which describes a huge fraction of
the corpus; the obvious relevant papers Doc 550 (\emph{laminar heat
transfer in tubes under slip-flow conditions}) and Doc 21 (\emph{on heat
transfer in slip flow}) are missed.

\textbf{Part--5 comparison.} All four systems are essentially unable to
fix Mode 3. On Q207 first-relevant ranks are 59 (VSM), 35 (BM25), 34
(LSA), 156 (QE) --- QE actively worsens. On Q117 only QE manages
P@10 = 0.10. On Q9 all four give P@10 = 0.00. The reason is structural:
BM25, LSA, and QE all operate on the surviving non-OOV tokens. None
introduces the missing semantic content. WordNet QE broadens generic
surviving terms with synonyms, which tends to amplify noise rather than
recover the missing technical compound. A real fix requires subword
tokenization or an embedding representation that does not depend on
exact-token vocabulary.

\section{Comparison with Part--5 systems}

Table~\ref{tab:partfive} reports aggregate metrics at $k=10$ for the four
systems alongside the paired Wilcoxon signed-rank significance of the
per-query AP@10 difference against VSM.

\begin{table}[h]
\centering
\small
\begin{tabular}{lcccccc}
\toprule
System & P@10 & R@10 & F$_{0.5}$@10 & MAP@10 & nDCG@10 & Wilcoxon vs.\ VSM \\
\midrule
VSM (baseline) & 0.2836 & 0.4100 & 0.2892 & 0.3026 & 0.4636 & --- \\
BM25 & \textbf{0.2893} & \textbf{0.4168} & \textbf{0.2950} & \textbf{0.3330} & \textbf{0.4926} & $p = 3.9\!\times\!10^{-6}$ \\
LSA ($k=300$) & 0.2640 & 0.3776 & 0.2687 & 0.2818 & 0.4357 & $p = 3.0\!\times\!10^{-3}$ \\
WordNet QE & 0.2596 & 0.3812 & 0.2657 & 0.2837 & 0.4344 & $p = 3.2\!\times\!10^{-6}$ \\
\bottomrule
\end{tabular}
\caption{Aggregate performance at $k=10$. Bold marks the best value in
each column. Wilcoxon $p$-values are computed on the per-query AP@10
differences against the VSM baseline.}
\label{tab:partfive}
\end{table}

The aggregate picture is more nuanced than the failure-mode analysis
suggests. BM25 is the only system that improves on every aggregate metric;
its gains are statistically significant. LSA and WordNet QE both
\emph{worsen} aggregate P@10 against the baseline despite significantly
different per-query distributions. This is consistent with the per-query
data: LSA helps strongly on synonymy queries (Mode 1) and slightly on
focus-drift queries (Mode 2), but introduces noise on queries the VSM
already answers well, because the latent projection blurs surface
specificity. WordNet QE has the same problem more acutely: synonym
expansion adds tokens that broaden recall but more often dilute precision
than recover missing concepts.

\section{Summary}

The TF--IDF VSM baseline achieves P@10 = 0.284 and nDCG@10 = 0.464 on
Cranfield. Failure-mode analysis distinguishes three structurally
different sources of error: synonymy (Mode 1), term-importance
limitations including negation and focus drift (Mode 2), and OOV tokens
(Mode 3). The Part--5 systems improve the baseline unevenly: BM25
delivers a consistent significant improvement across all aggregate
metrics; LSA addresses Mode 1 but worsens aggregates; WordNet QE
behaves similarly to LSA in aggregate and helps essentially none of the
three failure modes consistently. None of the three Part--5 systems
addresses Mode 2 (which requires syntactic role modelling) or Mode 3
(which requires subword or embedding representations), bounding what
this class of improvements can deliver on top of the bag-of-words
baseline.

\end{document}
